\section{Parent cyclic space and endpoint concentration}

\subsection{Cyclic accessibility}
\label{app:parent-accessibility}

We prove Proposition~\ref{prop:parent-cyclic-accessibility}.  Because
$\mathfrak A_{\rm par}$ is a unital star-algebra,
$\mathfrak A_{\rm par}\ket{J_{n,k}}$ is invariant under every element of
the algebra and its adjoint, hence is reducing.  Every $H_a$ commutes with
$D_n$, and $\ket{J_{n,k}}$ is $D_n$-fixed, proving the first item.

All rates in Eq.~\eqref{eq:johnson-metropolis} are positive on the connected
Johnson graph.  The zero eigenvalue of $H_a$ is consequently simple, with
ground vector Eq.~\eqref{eq:cycle-parent-ground}.  In finite dimension its
spectral projector $P_a=\ket{\psi_a}\!\bra{\psi_a}$ is a polynomial in
$H_a$ and belongs to $\mathfrak A_{\rm par}$.  Positivity gives
$\langle\psi_a|J_{n,k}\rangle>0$, so
$P_a\ket{J_{n,k}}$ is a nonzero multiple of $\ket{\psi_a}$.  Finally, both
restrictions have this unique ground state, while
$\mathcal K_J^{\rm par}\subseteq\mathcal H^{D_n}$.  Rayleigh--Ritz
minimization over the smaller orthogonal complement proves
Eq.~\eqref{eq:parent-cyclic-fixed-comparison}.

\subsection{Endpoint concentration}
\label{app:parent-concentration}

At energy $u$, occupied and vacant vertices decompose into the same number
$k-u$ of positive cyclic runs.  Choosing the marked start and the two run
compositions gives
\begin{equation}
 N_u=\frac{n}{k-u}\binom{k-1}{u}^2,
 \qquad0\leq u<k,
\end{equation}
which is Eq.~\eqref{eq:cycle-shell-count}.  At $a_f=7$, the $u=1$
contribution is $N_1n^{-7}=n(k-1)n^{-7}=O(n^{-5})$.  Consecutive terms
$N_un^{-7u}$ have ratio at most $2k^2n^{-7}=O(n^{-5})$.  Summing the
resulting geometric tail proves
Eq.~\eqref{eq:cycle-endpoint-concentration}.

\subsection{Adiabatic derivative bounds}
\label{app:parent-adiabatic}

A Johnson move removes one selected cycle vertex and inserts one unselected
vertex, so $|U(T)-U(S)|\leq2$, and each subset has
$k(n-k)=k^2$ neighbours.  Differentiating the matrix elements in
Eqs.~\eqref{eq:johnson-metropolis} and \eqref{eq:cycle-parent}, then bounding
the operator norm by the maximum absolute row sum, gives
\begin{equation}
 \|\partial_sH_{7s}\|\leq c_1n\log n,
 \qquad
 \|\partial_s^2H_{7s}\|\leq c_2n(\log n)^2,
 \label{eq:cycle-parent-derivative-bounds}
\end{equation}
with constants independent of $n$.
Proposition~\ref{prop:parent-cyclic-accessibility} and
Eq.~\eqref{eq:uniform-accessible-gap} give a simple ground state and minimum
cyclic gap at least $1024n^{-13}$.

Apply the quantitative gapped adiabatic bound
\cite{JansenRuskaiSeiler2007} to
$H_{7s}|_{\mathcal K_J^{\rm par}}$; restriction cannot increase either
derivative norm.  The three contributions scale as
$n^{27}\log n$, $n^{27}(\log n)^2$, and
$n^{41}(\log n)^2$, divided by $T$.  The last dominates and proves
Eq.~\eqref{eq:cycle-adiabatic-runtime}.  The final measurement statement
follows from Eq.~\eqref{eq:cycle-endpoint-concentration} and continuity of
measurement probabilities in state distance.
